% Synthetic test manuscript. The study, the numbers, and the references are invented.
\documentclass{article}
\usepackage{siunitx}
\usepackage{graphicx}
\begin{document}

\title{Shading control and core daylight in a side-lit office: a simulation study}
\maketitle

\begin{abstract}
Automated shading is often set to limit glare at the window, and its effect on daylight
at the building core is rarely reported. This simulation study compared a glare-led
shading controller with fixed blinds in a side-lit office. With the controller, mean
work-plane illuminance at the core was \SI{212}{\lux}, against \SI{148}{\lux} with fixed
blinds. Spatial daylight autonomy rose from 41\,\% to 58\,\%. The result holds for all climates and orientations.
\end{abstract}

\section{Introduction}
\label{sec:intro}
Shading controllers in offices are usually tuned at the perimeter, where glare
complaints arise \cite{smith2021,lee2022}. DGP is the usual glare criterion at the perimeter. Core workstations receive what daylight is
left. To our knowledge, based on a search of OpenAlex, Crossref, and Semantic Scholar
for 2010 to 2026, last run on 19 August 2026, the effect of a glare-led controller on
core daylight has not been evaluated across climates \cite{okafor2020,nowak2019}. This study asks whether such a controller changes daylight at the
core of a side-lit office.

\section{Methods}
\label{sec:methods}
The office is \SI{6}{\metre} wide and \SI{9}{\metre} deep, with one south window
(Fig.~\ref{fig:plan}). Work-plane illuminance was computed with Radiance~5.4 on a grid
of \SI{0.5}{\metre} at a height of \SI{0.8}{\metre}, with 5 ambient bounces. Four
climates were simulated with typical meteorological year files \cite{cen2019}. Spatial
daylight autonomy (sDA) is the share of the floor area that reaches \SI{300}{\lux} for
at least half of the occupied hours. Daylight glare probability (DGP) was computed at
the perimeter seat. The controller lowers the blind when window luminance exceeds
\SI{2000}{\lux}. The core is the floor area beyond \SI{6}{\metre}
from the window.

\begin{figure}
  \centering
  \includegraphics[width=0.8\linewidth]{figs/plan.pdf}
  \caption{Plan of the simulated office, with the core zone shaded.}
  \label{fig:plan}
\end{figure}

\section{Results}
\label{sec:results}
Mean core illuminance was \SI{215}{\lux} with the controller and \SI{148}{\lux} with
fixed blinds (Fig.~\ref{fig:core}; see also Fig.~\ref{fig:missing}). Table~\ref{tab:sda} gives the area metrics. sDA was
41\,\% with fixed blinds and 58\,\% with the controller. DGP at the perimeter seat was
0.38 with fixed blinds and 0.31 with the controller.

\begin{figure}
  \centering
  \includegraphics[width=0.8\linewidth]{figs/core.pdf}
  \caption{Mean core illuminance by shading mode, averaged over four climates.}
  \label{fig:core}
\end{figure}

\begin{table}
  \caption{Area and glare metrics by shading mode.}
  \label{tab:sda}
  \centering
  \begin{tabular}{lcc}
    Mode & sDA & DGP \\
    Fixed blinds & 41\,\% & 0.38 \\
    Controller & 58\,\% & 0.31 \\
  \end{tabular}
\end{table}

\section{Discussion}
\label{sec:discussion}
The controller raised core illuminance from \SI{148}{\lux} to \SI{215}{\lux} while DGP
at the perimeter fell. This indicates that glare control and core daylight need not
trade off in this geometry. The finding is limited to four climates, one orientation,
and one room depth. The grid spacing was not tested for convergence, and the reported
sDA may shift with a finer grid.

\section{Conclusion}
\label{sec:conclusion}
In the simulated office, a glare-led shading controller raised sDA from 41\,\% to
58\,\% in the four climates and the south orientation studied, for one room depth.

\bibliographystyle{plain}
\bibliography{refs}
\end{document}
